% TOP_PROBLEM: {"id": 7500112, "title": "Cancellation and Schröder–Bernstein for Torsion Abelian Groups", "category_id": 18, "category": {"id": 18, "name": "logic", "display_name": "Logic", "description": "Problems in mathematical logic, model theory, proof theory, and finite model theory.", "slug": "logic", "order_index": 18, "created_at": "2026-07-31T15:26:25.671Z"}, "status": "open"}
% FIRST_POSTED: null
% ENTRY_KIND: statement_only
% Imported from: batch20.tex; UnsolvedMath ID 7500112
\documentclass[11pt]{article}
\usepackage[a4paper,margin=25mm]{geometry}
\usepackage{amsmath,amssymb,mathtools}
\usepackage{fontspec,unicode-math}
\setmainfont{Latin Modern Roman}
\setsansfont{Latin Modern Sans}
\setmathfont{Latin Modern Math}
\usepackage{xurl,hyperref}
\hypersetup{colorlinks=true,linkcolor=blue,urlcolor=blue,pdftitle={UnsolvedMath Problem 7500112}}
\newcommand{\sref}[2]{\hyperlink{source:#1:#2}{[S#2]}}
\newcommand{\cataloguescope}[1]{\paragraph{Mathematical question.}#1}
\begin{document}
% UnsolvedMath ID 7500112; source catalogue code AMR-074-0112
\setcounter{section}{7500111}
\section{Cancellation and Schröder–Bernstein for Torsion Abelian Groups}\label{7500112}
{\small\sffamily UnsolvedMath ID 7500112 · Logic}
\subsection{Definitions and mathematical statement}

\paragraph{Shared notation.}$\mathbb N=\{1,2,\ldots\}$, and $\mathbb Z,\mathbb Q,\mathbb R,\mathbb C$ denote the integers, rationals, reals and complex numbers. Any other conventions are specified locally.


Work in second-order arithmetic over $\mathsf{RCA}_0$, the base theory with $\Delta^0_1$ comprehension and $\Sigma^0_1$ induction. The strength of a principle means which subsystems prove it and which subsystems it implies over that base; equivalence requires both implications.  The theory $\Pi^1_1\text{-}\mathsf{CA}_0$ adds comprehension for formulas with a universal block of set quantifiers and an arithmetical matrix. The theory $\mathsf{ATR}_0$ permits arithmetical transfinite recursion along every coded countable well-order. A countable abelian group is torsion if every element has finite order. A direct summand $G$ of $H$ means $H\cong G\oplus K$ for some group $K$. A group is divisible if for every positive integer $n$ and element $a$ there is $b$ with $nb=a$; it is reduced if it has no nonzero divisible subgroup.
\paragraph{Statement recorded in the supplied source.}
Is each of the following principles equivalent to $\Pi^1_1\text{-}\mathsf{CA}_0$? (i) For all countable torsion abelian $G,H$, $G\oplus G\cong H\oplus H$ implies $G\cong H$. (ii) For all such $G,H$, if each is a direct summand of the other, then $G\cong H$. In the source’s accompanying conjecture, is each principle restricted to reduced torsion abelian groups equivalent to $\mathsf{ATR}_0$? \sref{7500112}{2}
\subsection{Short English statement}
Do the two cancellation principles for countable torsion abelian groups each require analytic comprehension? When the groups have no divisible part, do they instead each require arithmetic transfinite recursion?
\subsection{Sources}
\begin{description}
\item[\hypertarget{source:7500112:1}{S1}] ULAM.AI and the UnsolvedMath Contributors, UnsolvedMath: A Curated Collection of Open Mathematics Problems, record 7500112 (AMR-074-0112). CC BY 4.0. \url{https://huggingface.co/datasets/ulamai/UnsolvedMath}.
\item[\hypertarget{source:7500112:2}{S2}] Antonio Montalbán, Open Questions in Reverse Mathematics, author preprint dated 17 August 2010, Question 12, printed p. 6. \url{https://math.berkeley.edu/~antonio/papers/questionsRM.pdf}.
\end{description}
\label{end:7500112}
\end{document}
